\chapter{Obligations d'entreprise}\label{ch:m2:corporate-bonds}

Les obligations d'une entreprise sont notées BBB$-$, le plus bas échelon de l'\omterm{def:m2:corporate-bonds:rating}{investment grade}. Depuis
des mois, le spread de ses obligations par rapport aux Treasuries ne cesse de s'élargir, à mesure que les analystes et
les investisseurs lisent les mêmes chiffres que les agences de notation. Puis une agence la dégrade
à BB$+$. Les obligations vont quitter les indices \omterm{def:m2:corporate-bonds:rating}{investment grade}~; les fonds dont
le mandat n'autorise que l'\omterm{def:m2:corporate-bonds:rating}{investment grade} sont tenus de les vendre, et les fonds \omterm{def:m2:corporate-bonds:rating}{high yield}, qui ne les détenaient pas la veille, doivent décider à quel prix
acheter. Lorsque la Banque centrale européenne a étudié ces «~\omterm{def:m2:corporate-bonds:rating}{fallen angels}~», elle a constaté
que l'essentiel de la baisse de leurs cours s'était produit avant la dégradation, et que
après la dégradation par la dernière agence, leurs cours se sont en partie redressés, comme si
les ventes forcées les avaient poussés trop loin. Ce chapitre présente les
obligations d'entreprise~: leur émission, la signification de leurs notations, les différents spreads
par lesquels leurs cours sont cotés et comparés, ainsi que les \omterm{def:m2:corporate-bonds:covenant}{covenants} et le \omterm{def:m2:corporate-bonds:covenant}{rang de séniorité} qui
déterminent ce que recouvrent leurs détenteurs lorsqu'une entreprise fait défaut.

\section{Émission et prime d'émission}

\begin{definition}[Prime d'émission]\label{def:m2:corporate-bonds:nic}
La \emph{prime d'émission}\index{prime d'émission} est le spread supplémentaire
auquel la nouvelle obligation d'une entreprise est cotée au-dessus de ses obligations existantes
de même échéance, le prix que l'émetteur paie pour attirer assez d'acheteurs rapidement.
\end{definition}

Une entreprise émet ses obligations par l'intermédiaire d'un syndicat de banques qui en assurent la prise ferme (voir
le livre~1, chapitre~2)~: elles annoncent l'opération, recueillent les ordres des investisseurs pendant
quelques heures, à partir d'une indication de prix initiale, resserrent le prix à mesure que le carnet
d'ordres s'étoffe, puis procèdent à l'allocation. Les obligations sont cotées en spread par rapport
à un Treasury de référence de même échéance. Lorsqu'une nouvelle émission se cote plus large que les
obligations en circulation de l'entreprise, l'écart constitue la prime~: une rémunération pour l'absorption d'une offre nouvelle
et pour l'engagement de montants importants en une seule fois. Une prime qui disparaît dès les premiers
jours de négociation est un gain pour les investisseurs et un coût pour l'émetteur.

\section{Notations et frontière de l'investment grade}

\begin{definition}[Notation de crédit, investment grade, high yield, fallen angel]\label{def:m2:corporate-bonds:rating}
Une \emph{notation de crédit}\index{notation de crédit} est l'avis d'une agence de notation sur
la solvabilité d'un émetteur ou d'une obligation, selon une échelle de lettres allant de AAA (Aaa) à D.
Les obligations notées BBB$-$ (Baa3) ou au-dessus sont \emph{investment grade}\index{investment grade}~;
celles notées BB$+$ (Ba1) ou en dessous sont \emph{high yield}\index{high yield}. Un
\emph{fallen angel}\index{fallen angel} est un émetteur dont la notation a été abaissée de l'investment
grade au high yield par au moins une des trois grandes agences.
\end{definition}

La frontière entre BBB$-$ et BB$+$ n'est qu'un cran sur une échelle, mais elle partage le
marché. De nombreux assureurs, fonds de pension et fonds obligataires ne peuvent détenir que de l'\omterm{def:m2:corporate-bonds:rating}{investment
grade}~; les indices sont scindés à cette frontière, de même que les fonds qui les répliquent~;
les régulateurs et les banques centrales s'en servent aussi. Lorsque la Réserve fédérale a mis en place
son dispositif d'achat d'obligations d'entreprise nouvellement émises en 2020, elle exigeait que les émetteurs soient
notés au moins BBB$-$/Baa3 au 22 mars 2020, et elle admettait ceux déclassés par la suite
jusqu'à BB$-$/Ba3~: la frontière, assortie d'un pont pour les \omterm{def:m2:corporate-bonds:rating}{fallen angels}.

L'étude de la BCE décrit le comportement du marché autour de cette frontière~: les primes de
credit default swap (CDS) des \omterm{def:m2:corporate-bonds:rating}{fallen angels} se sont élargies bien avant la première dégradation, ont peu bougé
le jour même, et se sont en partie redressées après la perte de la dernière notation \omterm{def:m2:corporate-bonds:rating}{investment grade}~;
les fonds qui répliquent les indices recourent surtout à l'échantillonnage et peuvent garder une obligation déclassée pendant
quelque temps, ce qui étale les ventes dans le temps. Parmi les \omterm{def:m2:corporate-bonds:rating}{fallen angels} recensés par Moody's sur vingt ans,
près d'un quart sont revenus en \omterm{def:m2:corporate-bonds:rating}{investment grade}, près de la moitié
sont restés en \omterm{def:m2:corporate-bonds:rating}{high yield} et 12~\% ont fait défaut. La \cref{fig:m2:corporate-bonds:angel}
illustre ce comportement sur une obligation d'illustration.

\begin{omfigure}
\begin{tikzpicture}
\begin{axis}[width=11cm,height=4.8cm,
  xlabel={mois écoulés depuis la dégradation qui fait sortir l'obligation de l'investment grade},ylabel={spread (\%)},
  tick label style={font=\small},label style={font=\small},
  xmin=-6.5,xmax=6.5,ymin=2.3,ymax=4.3,grid=major,grid style={gray!25},
  table/col sep=comma]
\addplot[omThm,very thick,mark=*,mark size=1.6pt] table[x=month,y=spread]{figdata/markets-2/21-corporate-bonds/angel.csv};
\draw[omDef,dashed] (axis cs:-6.5,3.5) -- (axis cs:6.5,3.5);
\node[omDef,font=\scriptsize,anchor=south west] at (axis cs:-6.3,3.52) {juste valeur pour BB$+$};
\draw[gray,dashed] (axis cs:0,2.3) -- (axis cs:0,4.3);
\node[omProp,font=\scriptsize,anchor=west] at (axis cs:0.15,4.08) {vente forcée};
\end{axis}
\end{tikzpicture}

\omcaption{Le spread d'un \omterm{def:m2:corporate-bonds:rating}{fallen angel} fictif~: l'essentiel de l'élargissement intervient
avant la dégradation, un dépassement dû aux ventes forcées le jour même, puis une remontée vers
le niveau juste pour sa nouvelle notation, le comportement observé par la BCE dans les credit default
swaps des \omterm{def:m2:corporate-bonds:rating}{fallen angels}. Trajectoire fictive~; données~: le
tutoriel du chapitre.}\label{fig:m2:corporate-bonds:angel}
\end{omfigure}

\section{La famille des mesures de spread}

\begin{definition}[Spread de crédit~; spreads G-, I-, Z- et d'asset-swap]\label{def:m2:corporate-bonds:spreads}
Un \emph{spread de crédit}\index{spread de crédit} est le rendement supplémentaire qu'une obligation risquée verse par
rapport à une courbe de référence. Le \emph{G-spread}\index{G-spread} est le rendement de l'obligation moins le
rendement souverain interpolé à son échéance~; le \emph{I-spread}\index{I-spread}
est son rendement moins le taux de swap interpolé. Le \emph{Z-spread}\index{Z-spread} est
la marge constante qui, ajoutée à la courbe zéro-coupon des swaps, actualise les flux de trésorerie
de l'obligation jusqu'à son prix. Le \emph{spread d'asset swap}\index{spread d'asset swap} est
la marge au-dessus du taux variable que gagne un acheteur de l'obligation en échangeant ses coupons fixes
contre des paiements variables, au pair~: l'écart entre la valeur de l'obligation sur la courbe de swap
et son prix, divisé par l'annuité de la jambe variable.
\end{definition}

Les quatre mesures répondent à des questions différentes. Le \omterm{def:m2:corporate-bonds:spreads}{G-spread} compare l'obligation avec
l'actif le plus sûr~; il mêle le crédit, le \omterm{def:m2:interest-rate-swaps:spread}{spread de swap} et tout écart entre le coupon et la
forme de la courbe. Le \omterm{def:m2:corporate-bonds:spreads}{I-spread} la compare aux swaps, les instruments auxquels les dealers recourent pour
se couvrir. Le \omterm{def:m2:corporate-bonds:spreads}{Z-spread} utilise toute la courbe zéro-coupon, et non un seul point, et valorise donc
les flux réels de l'obligation. Le \omterm{def:m2:corporate-bonds:spreads}{spread d'asset swap} est ce que gagne réellement un investisseur qui
se finance au taux variable et achète l'obligation avec un swap. Sur une courbe plate
et une obligation au pair, elles coïncident~; en général, il n'en est pas ainsi.

\begin{proposition}[Pourquoi les mesures diffèrent]\label{prop:m2:corporate-bonds:differ}
Pour une obligation de rendement $y$ et d'échéance $T$~: les spreads G et I diffèrent
du \omterm{def:m2:interest-rate-swaps:spread}{spread de swap} à $T$~; le \omterm{def:m2:corporate-bonds:spreads}{Z-spread} diffère du \omterm{def:m2:corporate-bonds:spreads}{I-spread} par l'effet de
la pente de la courbe sur les flux de l'obligation, effet qu'un rendement unique néglige~; et le
\omterm{def:m2:corporate-bonds:spreads}{spread d'asset swap} et le \omterm{def:m2:corporate-bonds:spreads}{Z-spread} mesurent le même écart entre le prix de l'obligation
et sa valeur sur la courbe de swap, le premier comme une marge courante versée au pair, le
second comme un décalage de tous les taux d'actualisation~: ils concordent au premier ordre au voisinage du pair, et
se séparent à mesure que le spread s'élargit et que le prix s'éloigne du pair.
\end{proposition}

\begin{proof}
Le premier point découle des définitions. Un rendement actualise chaque flux au même
taux~; le \omterm{def:m2:corporate-bonds:spreads}{Z-spread} actualise chacun au taux zéro de sa date, augmenté d'une constante~;
lorsque la courbe est pentue, les deux coïncident seulement si les taux zéro sont égaux au rendement.
L'asset swap transforme l'écart $P_{\mathrm{curve}} - P$ en une marge courante en le divisant par
l'annuité $A$ de la jambe variable. Un petit \omterm{def:m2:corporate-bonds:spreads}{Z-spread} $z$ fait baisser le prix
d'environ $z\sum_i t_i c_i P(0,t_i)$, une somme pondérée par la duration voisine de $100A$ pour une
obligation proche du pair~; les deux spreads sont donc proches. L'effet du \omterm{def:m2:corporate-bonds:spreads}{Z-spread} sur le prix n'est pas
linéaire, et les pondérations diffèrent lorsque le coupon est éloigné du taux de swap, de sorte que l'écart
s'élargit avec le spread et avec la distance au pair.
\end{proof}

\begin{example}[Une obligation, quatre spreads]\label{ex:m2:corporate-bonds:four}
Une obligation d'entreprise à sept ans verse un coupon semestriel de 5.50~\% et se négocie à 99.00~:
son rendement est de 5.6751~\%. Avec les courbes illustratives de
la \cref{fig:m2:corporate-bonds:curves}, le Treasury à sept ans à 4.50~\% et le swap à
4.25~\%, son \omterm{def:m2:corporate-bonds:spreads}{G-spread} vaut 117.5 points de base et son \omterm{def:m2:corporate-bonds:spreads}{I-spread} 142.5. Son \omterm{def:m2:corporate-bonds:spreads}{Z-spread} sur
la courbe des taux zéro de swap vaut 140.7 points de base et son \omterm{def:m2:corporate-bonds:spreads}{spread d'asset swap} 142.9.
\end{example}

\begin{omfigure}
\begin{tikzpicture}
\begin{axis}[width=11cm,height=5cm,
  xlabel={échéance (années)},ylabel={rendement (\%)},
  tick label style={font=\small},label style={font=\small},
  xmin=0,xmax=31,ymin=3.9,ymax=5.9,grid=major,grid style={gray!25},
  legend columns=3,legend style={font=\footnotesize,at={(0.5,-0.3)},anchor=north,draw=none,fill=none,/tikz/every even column/.append style={column sep=8pt}},
  table/col sep=comma]
\addplot[omDef,very thick] table[x=t,y=govt]{figdata/markets-2/21-corporate-bonds/curves.csv};
\addplot[omThm,very thick,dashed] table[x=t,y=swap]{figdata/markets-2/21-corporate-bonds/curves.csv};
\addplot[omProp,only marks,mark=*,mark size=2.4pt] table[x=t,y=y]{figdata/markets-2/21-corporate-bonds/bond.csv};
\draw[omProp,thick,{Stealth}-{Stealth}] (axis cs:7.6,4.50) -- (axis cs:7.6,5.675);
\node[omProp,font=\scriptsize,anchor=west] at (axis cs:7.9,5.05) {G~: 117.5 pb};
\draw[omThm,thick,{Stealth}-{Stealth}] (axis cs:6.4,4.25) -- (axis cs:6.4,5.675);
\node[omThm,font=\scriptsize,anchor=east] at (axis cs:6.1,4.85) {I~: 142.5 pb};
\legend{courbe au pair du Treasury,courbe au pair du swap,l'obligation}
\end{axis}
\end{tikzpicture}

\omcaption{L'obligation de la \cref{ex:m2:corporate-bonds:four} face aux courbes illustratives
du Treasury et du swap. Le \omterm{def:m2:corporate-bonds:spreads}{G-spread} est mesuré par rapport à la courbe du Treasury et le
\omterm{def:m2:corporate-bonds:spreads}{I-spread} par rapport à la courbe du swap, qui se situe ici sous celle-ci, comme elle le fait aux longues échéances
depuis que les \omterm{def:m2:interest-rate-swaps:spread}{spreads de swap} sont devenus négatifs (\cref{ch:m2:interest-rate-swaps}).
Illustratif~; données~: tutoriel du chapitre.}\label{fig:m2:corporate-bonds:curves}
\end{omfigure}

\begin{definition}[Spread ajusté des options]\label{def:m2:corporate-bonds:oas}
Le \emph{spread ajusté des options}\index{spread ajusté des options} (option-adjusted spread, OAS) d'une obligation dotée
d'options incorporées est le spread sur la courbe auquel la valeur de l'obligation, calculée
avec un modèle de taux qui valorise les options, égale son prix~: le \omterm{def:m2:corporate-bonds:spreads}{Z-spread} diminué
du coût des options, exprimé en spread.
\end{definition}

Le \omterm{def:m2:corporate-bonds:spreads}{Z-spread} d'une \omterm{def:m2:the-rates-options-market:swaption}{obligation remboursable par anticipation} surestime ce que son détenteur reçoit au titre du crédit, car
une partie du rendement rémunère l'option de remboursement anticipé que le détenteur a vendue à l'émetteur
(\cref{ch:m2:the-rates-options-market}). Si l'obligation de la
\cref{ex:m2:corporate-bonds:four} était remboursable au pair dans trois ans, cette option
serait une \omterm{def:m2:the-rates-options-market:swaption}{swaption receveuse} sur les quatre années restantes, dont le taux d'exercice serait le coupon.
Valorisée à une volatilité normale de 100 points de base, elle vaut 4.80 points, 79.7
points de base par an sur la durée de vie de l'obligation~: un \omterm{def:m2:corporate-bonds:oas}{spread ajusté des options} d'environ 61
points de base, contre un \omterm{def:m2:corporate-bonds:spreads}{Z-spread} de 141.

L'histoire des \omterm{def:m2:corporate-bonds:spreads}{spreads de crédit} est une suite de crises. Le Trésor publie une
courbe de rendement d'obligations d'entreprises de haute qualité, construite à partir d'obligations notées AAA, AA et A~; son
taux à dix ans, rapporté au rendement du Treasury à dix ans, est passé de moins de 1 point de pourcentage en
début 2007 à 5.04 en octobre 2008, puis à environ 2.1 au printemps 2020
(\cref{fig:m2:corporate-bonds:hqm}).

\begin{omfigure}
\begin{tikzpicture}
\begin{axis}[width=11cm,height=4.8cm,
  ylabel={points de pourcentage},
  tick label style={font=\small},label style={font=\small},
  xmin=1984,xmax=2026.75,ymin=0,ymax=5.5,grid=major,grid style={gray!25},
  xtick={1985,1990,1995,2000,2005,2010,2015,2020,2025},xticklabel style={/pgf/number format/1000 sep={}},
  table/col sep=comma]
\addplot[omThm,thick] table[x=t,y=spread]{figdata/markets-2/21-corporate-bonds/hqm.csv};
\end{axis}
\end{tikzpicture}

\omcaption{Taux zéro à dix ans des obligations d'entreprises de haute qualité (courbe HQM du Trésor, obligations
notées AAA, AA ou A), moins le rendement du Treasury à dix ans, mensuel, de 1984 à août 2026.
La mesure mêle le crédit et l'écart entre un \omterm{def:m2:government-bonds:zero}{taux zéro-coupon} et un \omterm{def:m2:government-bonds:zero}{taux au pair},
mais ses sommets, 2008 et 2020, correspondent aux crises de crédit. Données~: série FRED HQMCB10YR
(Trésor américain) et GS10.}\label{fig:m2:corporate-bonds:hqm}
\end{omfigure}

\begin{dated}{2026-09}{Spreads actuels}\label{dat:m2:corporate-bonds:spreads}
Taux zéro à dix ans HQM des obligations d'entreprises à 5.58~\% et Treasury à dix ans à 4.68~\% en août 2026
(moyennes mensuelles)~: un écart de 0.90 point de pourcentage. Les indices de spreads largement cités
d'obligations d'entreprises sont publiés par des fournisseurs d'indices sous licence et ne sont pas
reproduits ici.
\end{dated}

\section{Covenants, rang de séniorité et recouvrement}

\begin{definition}[Covenant, rang de séniorité, taux de recouvrement]\label{def:m2:corporate-bonds:covenant}
Un \emph{covenant}\index{covenant} est un engagement inscrit dans les conditions d'une obligation, qui restreint ce que
l'émetteur peut faire (emprunter davantage, verser des dividendes, céder des actifs, fusionner) ou l'oblige
à maintenir une certaine condition. Le \emph{rang de séniorité}\index{rang de séniorité} d'une créance est son rang dans
l'ordre de paiement en cas d'insolvabilité~: garantie avant senior non garantie, avant subordonnée, avant
subordonnée, avant capitaux propres. Le \emph{taux de recouvrement}\index{taux de recouvrement} est la part
de la valeur nominale d'une obligation que ses détenteurs reçoivent après le défaut de l'émetteur.
\end{definition}

\begin{omfigure}
\centering
\begin{tikzpicture}[font=\small,x=1cm,y=0.42cm]
\fill[omDef!25] (0,0) rectangle (2.4,6); \node[font=\footnotesize,align=center] at (1.2,3) {valeur de l'entreprise\\au défaut\\600};
\fill[omThm!30] (4,0) rectangle (6.4,2); \node[font=\footnotesize] at (5.2,1) {garantie 200};
\fill[omProp!30] (4,2) rectangle (6.4,7); \node[font=\footnotesize,align=center] at (5.2,4.5) {senior\\non garantie 500};
\fill[gray!25] (4,7) rectangle (6.4,9); \node[font=\footnotesize] at (5.2,8) {subordonnée 200};
\draw[thick,dashed] (2.6,6) -- (7.2,6);
\node[font=\footnotesize,anchor=west,align=left] at (7.3,1) {remboursée intégralement~: 100~\%};
\node[font=\footnotesize,anchor=west,align=left] at (7.3,4.2) {part de 400 sur 500~: 80~\%};
\node[font=\footnotesize,anchor=west,align=left] at (7.3,8) {rien ne reste~: 0~\%};
\end{tikzpicture}

\omcaption{\omterm{def:m2:corporate-bonds:covenant}{Rang de séniorité} et recouvrement sur un défaut illustratif. Une entreprise valant 600 au
défaut doit 200 garanties, 500 senior non garanties et 200 subordonnées~: les
créanciers garantis sont remboursés intégralement, les créances senior non garanties se partagent les 400 restants, et les
créances subordonnées ne récupèrent rien. Schéma.}\label{fig:m2:corporate-bonds:seniority}
\end{omfigure}

Les obligations \omterm{def:m2:corporate-bonds:rating}{investment grade} comportent en général peu de \omterm{def:m2:corporate-bonds:covenant}{covenants}~; les obligations \omterm{def:m2:corporate-bonds:rating}{high yield} en comportent davantage,
car leurs détenteurs ont besoin d'une protection contre une entreprise qui s'endette davantage
ou qui verse sa trésorerie aux actionnaires. Le \omterm{def:m2:corporate-bonds:covenant}{rang de séniorité} détermine le recouvrement~: en cas de liquidation ou
de restructuration, les créances senior garanties sont payées en premier sur leur collatéral, les
créances senior non garanties se partagent ce qui reste, et les créances subordonnées ne sont payées
qu'après elles (\cref{fig:m2:corporate-bonds:seniority}). Les \omterm{def:m2:corporate-bonds:covenant}{taux de recouvrement} varient fortement selon l'entreprise, le secteur et le cycle, ce qui
explique qu'un spread soit une rémunération pour deux grandeurs incertaines~: la probabilité de
défaut et la perte subie lorsqu'il se produit, objet
du \cref{ch:m2:credit-default-swaps}.

\begin{remark}[Spread, défaut et perte]\label{rem:m2:corporate-bonds:loss}
Sur un an, un spread $s$ rémunère à peu près une probabilité annuelle de défaut $p$ multipliée par une
perte en cas de défaut $1 - R$~: $s \approx p(1 - R)$, auxquels s'ajoutent des primes pour le risque et
l'illiquidité. Un spread de 140 points de base avec un recouvrement de 40~\% correspond, si tout
était intégralement une perte attendue, à une probabilité de défaut de 2.3~\% par an. L'écart
entre le taux de défaut effectif des obligations et ce niveau est la prime que les investisseurs exigent pour
porter le risque de défaut et l'illiquidité.
\end{remark}

\section{Tutoriel~: quatre spreads d'une même obligation}\label{tut:m2:corporate-bonds:spreads}

\begin{tutorial}
\textbf{Objectif.} Calculer les \omterm{def:m2:corporate-bonds:spreads}{G-spread}, \omterm{def:m2:corporate-bonds:spreads}{I-spread}, \omterm{def:m2:corporate-bonds:spreads}{Z-spread} et \omterm{def:m2:corporate-bonds:spreads}{spread d'asset swap} d'une obligation, ainsi que son
\omterm{def:m2:corporate-bonds:oas}{spread ajusté des options} (OAS) si elle est remboursable par anticipation (callable)~; tracer l'historique du spread.
\textbf{État final~:}
les \cref{fig:m2:corporate-bonds:curves,fig:m2:corporate-bonds:hqm},
l'\cref{ex:m2:corporate-bonds:four} et les chiffres du problème du week-end.
\begin{tutsteps}
\item \textbf{Courbes et rendements}~: interpolation, flux de trésorerie de l'obligation, prix et
rendement.
\omcode{code/firm/spreads/firm_spreads.py}{10}{33}{Interpolation des courbes, flux de trésorerie, prix et rendement.}\label{lst:m2:corporate-bonds:yield}
\item \textbf{Les spreads}~: vers les courbes au pair, vers la courbe zéro, et l'asset swap.
\omcode{code/firm/spreads/firm_spreads.py}{36}{61}{G-spread, I-spread, Z-spread et spread d'asset swap.}\label{lst:m2:corporate-bonds:spreads}
\item \textbf{Exécuter} \texttt{credit\_demo.measures()},
\texttt{credit\_demo.callable\_oas()}, \texttt{credit\_demo.fallen\_angel()} et
\texttt{fig\_credit.py}.
\end{tutsteps}
\textbf{Pour aller plus loin.} Accentuer la pente de la courbe des swaps et observer le \omterm{def:m2:corporate-bonds:spreads}{Z-spread}
s'éloigner du \omterm{def:m2:corporate-bonds:spreads}{I-spread}~; fixer ensuite le prix de l'obligation à 110 puis à 90, et comparer les
\omterm{def:m2:corporate-bonds:spreads}{spreads d'asset swap} et \omterm{def:m2:corporate-bonds:spreads}{Z-spread}, comme le prédit la \cref{prop:m2:corporate-bonds:differ}.
\end{tutorial}

\section{Construction~: la bibliothèque de mesures de spread}\label{bld:m2:corporate-bonds:spreads}

\begin{build}
\bfield{Rôle} La firme miniature cote et compare des obligations d'entreprises, les couvre
avec des Treasuries ou des swaps, et doit parler la langue des spreads de chaque client.
\bfield{Interface} \texttt{interp(curve, t)}~; \texttt{flows}~; \texttt{price\_from\_yield}~;
\texttt{yield\_from\_price}~; \texttt{g\_spread}~; \texttt{i\_spread}~;
\texttt{price\_on\_zero}~; \texttt{z\_spread}~; \texttt{asset\_swap\_spread}~;
\texttt{spread\_price\_impact(duration, widening, price)}.
\bfield{Règles} Courbes définies comme des points (années, taux), interpolées linéairement~; coupons semestriels
prix coupon inclus à une date de coupon~; taux zéro en capitalisation continue~; recherche
de racine par bissection.
\bfield{Tests d'acceptation} \texttt{code/firm/spreads/tests/}~: allers-retours d'interpolation et
de rendement~; spreads égaux sur les courbes plates~; le \omterm{def:m2:corporate-bonds:spreads}{Z-spread} retrouvé à partir d'un prix
construit avec lui~; un \omterm{def:m2:corporate-bonds:spreads}{spread d'asset swap} nul pour une obligation évaluée sur la courbe.
\bfield{Pour aller plus loin} Coupon couru et dates de règlement-livraison (la construction
d'obligation du \cref{ch:m2:government-bonds})~; OAS pour les calls bermudéens avec un treillis de taux court
(livre~6)~; spreads par rapport à une courbe d'émetteur ajustée.
\end{build}

\begin{omsources}
\item M.~Belloni, T.~Helmersson, M.~Jarmuzek, B.~Mosk et F.~Nikolic,
«~Understanding what happens when angels fall~», ECB \emph{Financial Stability Review},
novembre 2020.
\item Federal Reserve, Primary Market Corporate Credit Facility term sheet, avril 2020.
\item US Department of the Treasury, High Quality Market corporate bond yield curve~;
séries FRED HQMCB10YR et GS10.
\end{omsources}

\section{Exercices}

\begin{exercise}[$\star$]\label{exo:m2:corporate-bonds:1}
Une obligation à dix ans rapporte 6.10~\%. Le Treasury à dix ans rapporte 4.68~\% et le
taux de swap à dix ans est de 4.40~\%. Donner ses \omterm{def:m2:corporate-bonds:spreads}{G-spread} et \omterm{def:m2:corporate-bonds:spreads}{I-spread}.
\end{exercise}

\begin{exercise}[$\star$]\label{exo:m2:corporate-bonds:2}
Notée Baa3 par une agence et BB$+$ par une autre, l'obligation est-elle \omterm{def:m2:corporate-bonds:rating}{investment grade}~?
Son émetteur est-il un \omterm{def:m2:corporate-bonds:rating}{fallen angel} s'il était noté BBB$-$ auprès des deux agences le mois dernier~?
\end{exercise}

\begin{exercise}[$\star$]\label{exo:m2:corporate-bonds:3}
Une obligation de \omterm{def:m2:government-bonds:duration}{duration modifiée} 6 voit son spread s'élargir de 80 points de base. Donner sa
variation de prix pour 100.
\end{exercise}

\begin{exercise}[$\star\star$]\label{exo:m2:corporate-bonds:4}
Dans l'\cref{ex:m2:corporate-bonds:four}, le \omterm{def:m2:corporate-bonds:spreads}{Z-spread} et le \omterm{def:m2:corporate-bonds:spreads}{spread d'asset swap} diffèrent
de 2.2 points de base. Pourquoi sont-ils proches, et quand pourraient-ils différer davantage~?
\end{exercise}

\begin{exercise}[$\star\star$]\label{exo:m2:corporate-bonds:5}
Un spread de 140 points de base et un recouvrement de 40~\%~: quelle probabilité annuelle de défaut
le spread rémunérerait-il s'il était entièrement une perte attendue~? Que rémunère-t-il
d'autre~?
\end{exercise}

\begin{exercise}[$\star\star$]\label{exo:m2:corporate-bonds:6}
Pourquoi une nouvelle émission s'accompagne-t-elle d'une \omterm{def:m2:corporate-bonds:nic}{prime d'émission}, et qui gagne si elle disparaît~?
\end{exercise}

\begin{exercise}[$\star\star\star$]\label{exo:m2:corporate-bonds:7}
\emph{Programmation.} Avec la construction, recalculer le prix de l'obligation de l'\cref{ex:m2:corporate-bonds:four}
à 105 et donner ses quatre spreads.
\end{exercise}

\begin{exercise}[$\star\star\star$]\label{exo:m2:corporate-bonds:8}
\emph{Repérez l'erreur.} «~Cette \omterm{def:m2:the-rates-options-market:swaption}{obligation remboursable par anticipation} (callable) rapporte 140 points de base au-dessus des swaps, 
comme l'obligation bullet du même émetteur~; elle est donc tout aussi bon marché.~» Corrigez-la.
\end{exercise}

\section{Problème~: Le fallen angel}

\begin{problem}[{Problème du week-end --- coût et gain d'une vente forcée}]\label{pb:m2:corporate-bonds:1}
Un fonds détient 50~millions de dollars d'obligations d'une entreprise, de \omterm{def:m2:government-bonds:duration}{duration modifiée} 6. Six mois
avant toute dégradation, leur spread est de 250 points de base~; il s'élargit à 290 trois mois
avant, à 330 un mois avant, et à 400 le jour où la dernière agence ramène la note de l'entreprise à
BB$+$ et où les obligations sortent des indices \omterm{def:m2:corporate-bonds:rating}{investment grade}~; puis 380 un mois plus tard, 350
après trois mois et 345 après six mois. Les analystes jugent 350 juste pour un crédit BB$+$.

\medskip\noindent\textbf{Partie I --- Avant la dégradation.}
\begin{enumerate}
\item Combien les obligations ont-elles perdu, pour 100, entre six mois avant et la dégradation~?
\item Quelle part de cette baisse s'est produite avant le jour de la dégradation~?
\item Pourquoi le marché bouge-t-il avant les agences~?
\item Qu'aurait fait un fonds \omterm{def:m2:corporate-bonds:rating}{investment grade} face à une perspective négative~?
\item Donnez la réaction du jour, pour 100.
\end{enumerate}

\medskip\noindent\textbf{Partie II --- La vente forcée.}
\begin{enumerate}[resume]
\item Donnez le dépassement au-delà de la juste valeur, en spread et pour 100.
\item Que coûte ce dépassement au fonds s'il vend le jour même~?
\item Qui achète, et pourquoi a-t-il besoin d'une décote~?
\item Comment les règles des indices et l'échantillonnage du fonds modifient-ils la pression~?
\item Que gagne le fonds à attendre trois mois, s'il en a la possibilité~?
\end{enumerate}

\medskip\noindent\textbf{Partie III --- L'acheteur.}
\begin{enumerate}[resume]
\item Donnez le gain d'un acheteur, pour 100, à acheter le jour même et à revendre trois mois plus tard.
\item Quel risque l'acheteur supporte-t-il~?
\item Que disent les chiffres de la BCE sur les résultats ultérieurs au sujet de ce risque~?
\item Comment le carry du spread plus élevé s'ajoute-t-il au rendement~?
\item Pourquoi la transaction n'est-elle pas un gain sans risque~?
\end{enumerate}

\medskip\noindent\textbf{Partie IV --- Jugement.}
\begin{enumerate}[resume]
\item Pourquoi les notations comptent-elles autant lorsque les prix bougent en premier~?
\item Comment une facilité de banque centrale comme celle de 2020 modifierait-elle la configuration~?
\item Que demanderiez-vous que le mandat du fonds autorise~?
\item Énoncez le \emph{résultat nommé}~: la décote de vente forcée et sa résorption.
\item En une phrase~: pourquoi les \omterm{def:m2:corporate-bonds:rating}{fallen angels} dépassent-ils la juste valeur~?
\end{enumerate}
\end{problem}

\section{Questions d'entretien}

\begin{interviewq}[$\star$ \iqroles{trader, chercheur}]\label{iq:m2:corporate-bonds:1}
Quelle est la différence entre le \omterm{def:m2:corporate-bonds:spreads}{G-spread}, le \omterm{def:m2:corporate-bonds:spreads}{Z-spread} et le \omterm{def:m2:corporate-bonds:spreads}{spread d'asset swap}~?
\end{interviewq}

\begin{interviewq}[$\star$ \iqroles{trader, banque}]\label{iq:m2:corporate-bonds:2}
Qu'est-ce qu'un \omterm{def:m2:corporate-bonds:rating}{fallen angel}, et pourquoi ses obligations se déprécient-elles souvent autour de la dégradation~?
\end{interviewq}

\begin{interviewq}[$\star\star$ \iqroles{chercheur}]\label{iq:m2:corporate-bonds:3}
Comment décomposeriez-vous un \omterm{def:m2:corporate-bonds:spreads}{spread de crédit} en perte attendue et prime de risque~?
\end{interviewq}

\begin{interviewq}[$\star\star$ \iqroles{trader}]\label{iq:m2:corporate-bonds:4}
Un client demande un prix pour 50~millions de dollars d'une obligation BBB qui se traite rarement. Comment
la pricez-vous~?
\end{interviewq}

\begin{interviewq}[$\star\star$ \iqroles{chercheur, trader}]\label{iq:m2:corporate-bonds:5}
Pourquoi l'OAS d'une \omterm{def:m2:the-rates-options-market:swaption}{obligation remboursable par anticipation} diffère-t-il de son \omterm{def:m2:corporate-bonds:spreads}{Z-spread}, et lequel
faut-il comparer entre les obligations~?
\end{interviewq}

\begin{interviewq}[$\star\star\star$ \iqroles{développeur, chercheur}]\label{iq:m2:corporate-bonds:6}
Concevez un système qui calcule les spreads de cinquante mille obligations d'entreprises toutes les minutes.
\end{interviewq}
